\documentclass[10pt,twocolumn,a4paper]{article}
\usepackage[utf8]{inputenc}
\usepackage[T1]{fontenc}
\usepackage{amsmath,amssymb,amsthm}
\usepackage{gensymb}
\usepackage{mathtools}
\usepackage{physics}
\usepackage{geometry}
\usepackage{hyperref}
\usepackage{cleveref}
\usepackage{booktabs}
\usepackage{array}
\usepackage{multirow}
\usepackage{xcolor}
\usepackage{tikz}
\usepackage{pgfplots}
\pgfplotsset{compat=1.18}
\usepackage{enumitem}
\usepackage{float}
\usepackage{caption}
\usepackage{subcaption}
\usepackage{bm}

\geometry{margin=2.5cm}

\theoremstyle{definition}
\newtheorem{axiom}{Axiom}
\newtheorem{theorem}{Theorem}
\newtheorem{lemma}[theorem]{Lemma}
\newtheorem{corollary}[theorem]{Corollary}
\newtheorem{proposition}[theorem]{Proposition}
\newtheorem{definition}{Definition}
\newtheorem{remark}{Remark}

\newcommand{\SE}{S_{\text{e}}}
\newcommand{\SK}{S_{\text{k}}}
\newcommand{\ST}{S_{\text{t}}}
\newcommand{\Rc}{R_c}
\newcommand{\Rn}{R_n}
\newcommand{\kB}{k_{\text{B}}}
\newcommand{\kBT}{k_{\text{B}}T}
\newcommand{\PCHR}{\text{PCHR}}
\newcommand{\CSCI}{\text{CSCI}}
\newcommand{\RARS}{\text{RARS}}
\newcommand{\TCC}{\text{TCC}}

\title{\textbf{Sensor Disambiguation via Partition Coupling:\\
Deriving Multi-Modal Wearable Health Metrics\\from the Bounded Phase Space Framework}\\[0.5em]
}

\author{
Kundai Farai Sachikonye\\
AIMe Registry for Artificial Intelligence\\
\texttt{kundai.sachikonye@bitspark.com}
}

\date{\today}

\begin{document}

\maketitle

\begin{abstract}
Modern wearable devices measure heart rate, heart rate variability, blood oxygen
saturation, skin temperature, respiratory rate, and tri-axial acceleration ---
yet process each sensor stream independently, discarding the inter-scale coupling
that the Bounded Phase Space framework~\cite{sachikonye2024integration} proves is
the fundamental organising principle of physiology. We exploit this coupling to
derive five novel composite metrics from first principles: (1) the
\emph{Partition-Coupled Heart Rate} ($\PCHR$), which disambiguates identical
heart rate readings across distinct physiological regimes by incorporating
temperature and oxygen coupling; (2) the \emph{S-Entropy Health Coordinates}
$(\SK, \ST, \SE) \in [0,1]^3$, a three-dimensional health state that unifies
all sensor channels into a single navigable manifold; (3) the
\emph{Temperature-Corrected Coherence} ($\TCC$), which removes metabolic bias
from the Kuramoto order parameter $\Rc$; (4) the \emph{Cross-Scale Coherence
Index} ($\CSCI$), which quantifies how faithfully the multi-sensor coupling
matches partition-theoretic predictions; and (5) the \emph{Regime-Aware Recovery
Score} ($\RARS$), which tracks post-exercise recovery as a trajectory through
partition regime space rather than a scalar decay. We validate all five metrics
against 86 nights of sleep photoplethysmography and 104 days of activity data
from an Oura Ring, demonstrating that coupled metrics achieve superior
physiological discrimination compared to their single-sensor counterparts.
The framework yields a principled sensor fusion architecture applicable to
any wearable platform measuring two or more of the core physiological signals.
\end{abstract}

%============================================================
\section{Introduction}
\label{sec:introduction}
%============================================================

\subsection{The Disambiguation Problem}

A resting heart rate of 90 beats per minute conveys no physiological
specificity. The same reading arises during mild exercise recovery,
low-grade fever, acute psychological stress, early-stage heart failure,
altitude acclimatisation, and dehydration. Each condition occupies a
distinct region of the underlying physiological state space; the scalar
measurement $\text{HR} = 90$ is a degenerate projection that collapses
these distinct states onto a single number.

This degeneracy is not unique to heart rate. SpO$_2 = 94\%$ is normal
at 2500\,m altitude but pathological at sea level. RMSSD $= 45$\,ms
may indicate healthy vagal tone in an athlete or sympathetic withdrawal
in heart failure. Skin temperature deviation of +0.5°C may
reflect luteal phase hormonal shift, early infection, or post-exercise
thermogenesis. In every case, the single-sensor reading is
\emph{necessary but insufficient} for physiological classification.

Current wearable platforms address this limitation through heuristic
scoring algorithms (``readiness scores'', ``stress levels'') that
combine sensor channels via machine-learned weights without
physiological grounding~\cite{altini2021promise,derosario2019validity}.
These scores are opaque, non-transferable across populations, and
incapable of distinguishing novel physiological states not represented
in training data.

\subsection{The Partition-Theoretic Solution}

The Bounded Phase Space framework~\cite{sachikonye2024integration}
establishes that all physiological oscillators --- from mitochondrial
proton transfer ($10^{13}$\,Hz) through the cardiac cycle ($\sim$1\,Hz)
to circadian rhythms ($10^{-5}$\,Hz) --- share identical mathematical
structure as partition systems governed by the single axiom
$C(n) = 2n^2$~\cite{sachikonye2024fluid,sachikonye2024gas}. The
framework derives, from first principles, explicit coupling equations
between cardiac coherence ($\Rc$), neural coherence ($\Rn$), oxygen
tension ($P_{\text{O}_2}$), metabolic rate, and temperature --- the
same quantities that wearable sensors measure.

The key insight of this paper is that these coupling equations can be
\emph{inverted}: given simultaneous measurements from multiple sensors,
the partition coupling structure disambiguates the underlying
physiological state. Two identical HR readings that differ in their
temperature and SpO$_2$ context are mapped to different locations in
S-entropy coordinate space $(\SK, \ST, \SE)$, resolving the degeneracy
that plagues single-sensor interpretation.

\subsection{Contributions}

\begin{enumerate}[leftmargin=*]
\item We derive five novel composite health metrics from the partition
  coupling equations, each with an explicit functional form, physical
  units, and defined regime boundaries (Section~\ref{sec:metrics}).
\item We establish the formal conditions under which multi-sensor
  coupling resolves measurement degeneracy
  (Theorem~\ref{thm:disambiguation}).
\item We validate the metrics against 86 nights of sleep and 104 days
  of activity data, demonstrating improved physiological discrimination
  (Section~\ref{sec:validation}).
\item We provide an open sensor fusion architecture applicable to any
  wearable platform (Section~\ref{sec:architecture}).
\end{enumerate}

\subsection{Structure of the Paper}

Section~\ref{sec:framework} reviews the partition-theoretic foundations.
Section~\ref{sec:sensor_model} maps wearable sensor channels to
partition variables. Section~\ref{sec:metrics} derives the five
composite metrics. Section~\ref{sec:validation} presents empirical
validation. Section~\ref{sec:architecture} describes the sensor fusion
architecture. Section~\ref{sec:discussion} discusses clinical
implications, limitations, and future work.

%============================================================
\section{Partition-Theoretic Foundations}
\label{sec:framework}
%============================================================

\subsection{The Bounded Phase Space Axiom}

\begin{axiom}[Bounded Phase Space Law~\cite{sachikonye2024integration}]
\label{ax:bpsl}
Every physical system occupies a phase space of finite, integer-valued
capacity:
\begin{equation}
\boxed{C(n) = 2n^2}
\end{equation}
with partition coordinates $(n, \ell, m, s)$ satisfying
$n \geq 1$, $0 \leq \ell \leq n-1$, $-\ell \leq m \leq \ell$,
$s \in \{-\tfrac{1}{2}, +\tfrac{1}{2}\}$.
\end{axiom}

From this axiom, three entropy measures are proved equivalent~\cite{sachikonye2024integration}:
\begin{equation}
S_{\text{osc}} = S_{\text{cat}} = S_{\text{part}} = \kB \mathcal{M} \ln n
\end{equation}
establishing that oscillatory, categorical, and partition-theoretic
descriptions are interchangeable projections of the same underlying
structure.

\subsection{S-Entropy Coordinate Space}

\begin{definition}[S-Entropy Coordinates]
\label{def:sentropy}
The normalised information state of any biological oscillator is
described by:
\begin{equation}
\label{eq:sentropy}
(\SK, \ST, \SE) \in [0,1]^3
\end{equation}
where:
\begin{align}
\SK &= \frac{\mathcal{M} \ln n}{\mathcal{M}_{\max} \ln n_{\max}}
    \quad \text{(knowledge depth)} \label{eq:sk}\\
\ST &= 1 - e^{-t/\tau_{\text{char}}}
    \quad \text{(temporal integration)} \label{eq:st}\\
\SE &= \frac{S_{\text{actual}}}{S_{\max}}
    = \frac{\ln p(n, \mathcal{M})}{\ln p(n_{\max}, \mathcal{M}_{\max})}
    \quad \text{(entropy utilisation)} \label{eq:se}
\end{align}
\end{definition}

These coordinates are \emph{universal}: the same axes describe cardiac
ventricular states, neural ensemble firing patterns, metabolic substrate
allocation, and --- crucially for this paper --- the coupled multi-sensor
state of a wearable device.

\subsection{The Kuramoto Order Parameter and Regime Classification}

The Kuramoto order parameter~\cite{kuramoto1984chemical,strogatz2000kuramoto}
\begin{equation}
R e^{i\psi} = \frac{1}{N}\sum_{j=1}^{N} e^{i\theta_j}
\end{equation}
classifies both cardiac ($\Rc$) and neural ($\Rn$) dynamics into five
regimes (Table~\ref{tab:regimes}).

\begin{table}[H]
\centering
\caption{Partition Regime Classification}
\label{tab:regimes}
\begin{tabular}{lll}
\toprule
\textbf{Regime} & \textbf{$R$ range} & \textbf{Equation of state}\\
\midrule
Phase-locked & $>0.95$ & $S = 1 + K/\sigma(\omega)$\\
Coherent     & $0.80$--$0.95$ & $S = 1 + R^2/(1-R^2)$\\
Cascade      & $0.50$--$0.80$ & $S = \prod(1 + F_{\text{out}}/F_{\text{in}})$\\
Aperture     & $0.30$--$0.50$ & $S = n^2/n^2_{\max}$\\
Turbulent    & $<0.30$ & $S = 1 - \sigma^2(\phi)/(2\pi^2)$\\
\bottomrule
\end{tabular}
\end{table}

For wearable applications, $\Rc$ is estimated from beat-to-beat RR
intervals~\cite{malik1996heart,shaffer2017overview}:
\begin{equation}
\label{eq:rc_estimator}
\Rc = \exp\!\left(-2\pi^2 \cdot \text{CV}^2\right),
\quad \text{CV} = \frac{\text{RMSSD}}{\overline{T}_{RR}}
\end{equation}
where RMSSD is the root-mean-square of successive differences and
$\overline{T}_{RR} = 60{,}000/\text{HR}$ in milliseconds.

\subsection{Inter-Scale Coupling Equations}

The partition framework derives three coupling equations that connect
the quantities measured by wearable sensors:

\paragraph{Oxygen-neural coupling.}
\begin{equation}
\label{eq:kappa}
\frac{d\Rn}{dt} = \kappa_{\text{O}_2} \left(\Rn^{\text{target}}(P_{\text{O}_2}) - \Rn\right)
\end{equation}
with $\kappa_{\text{O}_2} = 4.7 \times 10^{-3}\,\text{s}^{-1}$~\cite{sachikonye2024integration}.

\paragraph{Cardiac-neural coupling.}
\begin{equation}
\label{eq:cn_coupling}
\frac{\Rn}{\Rc} = \frac{0.87}{\sqrt{\Rc}}
\quad \text{(coupled states)}
\end{equation}

\paragraph{Metabolic-thermal coupling.}
The metabolic rate $\dot{Q}_{\text{met}}$ couples to skin temperature
deviation $\Delta T_{\text{skin}}$ via the partition buffer
equation~\cite{sachikonye2024cardiovascular}:
\begin{equation}
\label{eq:thermal}
\Delta T_{\text{skin}}(t) = \frac{\dot{Q}_{\text{met}}}{h_c A_s}
\left(1 - e^{-t/\tau_{\text{therm}}}\right)
\end{equation}
where $h_c$ is the convective heat transfer coefficient, $A_s$ is skin
surface area, and $\tau_{\text{therm}} \approx 10$--$30$\,min is the
thermal time constant.

\begin{figure*}[!htbp]
\centering
\includegraphics[width=0.9\textwidth]{figures/panel19_csci_cross_scale.png}
\caption{\textbf{Cross-Scale Coherence Index: Detecting Inter-Scale Decoupling from Wrist Sensors (Oura Ring, 86 nights).}
\textbf{Panel A (left):} Violin plots of CSCI distribution per sleep stage. Deep sleep achieves the highest CSCI ($\sim$0.91), indicating strong adherence to partition-predicted coupling ratios between HR, RMSSD, and temperature. Awake shows moderate coupling ($\sim$0.88). REM exhibits the \emph{lowest} CSCI ($\sim$0.71)---independently quantifying the cardiac--neural active decoupling from wrist-worn cardiac sensors alone, without any neural (EEG) data. This is the first demonstration that inter-scale decoupling during REM is detectable purely from the departure of cardiac sensor coupling from partition-theoretic predictions.
\textbf{Panel B (center-left):} Scatter plot of CSCI versus Oura sleep score, with linear trend line. Higher CSCI (better inter-scale coupling) correlates positively with sleep quality, consistent with the framework prediction that physiological health corresponds to maintained coupling across the partition hierarchy.
\textbf{Panel C (center-right):} Three-dimensional scatter of $R_c$ $\times$ RMSSD $\times$ CSCI, colored by sleep stage. The three-dimensional structure reveals that CSCI provides information orthogonal to both $R_c$ and RMSSD: epochs with identical $R_c$ and RMSSD can have different CSCI values depending on whether the temperature coupling is maintained.
\textbf{Panel D (right):} Heatmap of coupling pair deviations per sleep stage, showing the fractional departure $|r^{\text{obs}} - r^{\text{pred}}| / |r^{\text{pred}}|$ for each of the sensor coupling pairs (HR--RMSSD, HR--temperature). REM shows the largest deviations across both pairs, while Deep sleep shows the smallest, confirming that the CSCI captures a genuine physiological signal rather than measurement noise.}
\label{fig:csci_cross_scale}
\end{figure*}

%============================================================
\section{Sensor-to-Partition Mapping}
\label{sec:sensor_model}
%============================================================

\subsection{Available Sensor Channels}

A modern wrist-worn wearable (Oura Ring, Apple Watch, Garmin, Whoop)
provides the following channels at the indicated temporal
resolutions~\cite{altini2021promise,derosario2019validity}:

\begin{table}[H]
\centering
\caption{Wearable Sensor Channels and Partition Variables}
\label{tab:sensors}
\begin{tabular}{llll}
\toprule
\textbf{Sensor} & \textbf{Signal} & \textbf{Resolution} & \textbf{Partition variable}\\
\midrule
PPG green & HR, IBI & 1\,s & $\Rc$, $\overline{T}_{RR}$\\
PPG green & RMSSD   & 5\,min & $\sigma^2(\phi)$, CV\\
PPG red/IR & SpO$_2$ & 1\,min & $P_{\text{O}_2}$\\
Thermistor & $T_{\text{skin}}$ & 1\,min & $\dot{Q}_{\text{met}}$\\
Accelerometer & Activity & 1\,min & MET, $\dot{W}_{\text{ext}}$\\
Derived & Resp.\ rate & 1\,min & Scale 8 oscillator\\
Derived & Hypnogram & 5\,min & Neural regime\\
\bottomrule
\end{tabular}
\end{table}

\subsection{The Degeneracy Problem: Formal Statement}

\begin{definition}[Sensor Degeneracy]
\label{def:degeneracy}
A scalar measurement $x_i$ from sensor $i$ is \emph{$k$-degenerate}
if there exist $k$ distinct physiological states
$\bm{s}_1, \ldots, \bm{s}_k \in \mathcal{M}$ such that
$\pi_i(\bm{s}_j) = x_i$ for all $j = 1, \ldots, k$, where
$\pi_i: \mathcal{M} \to \mathbb{R}$ is the projection onto sensor $i$.
\end{definition}

\begin{theorem}[Disambiguation by Partition Coupling]
\label{thm:disambiguation}
Let $\mathcal{M} = (R, \sigma^2, \SK, \ST, \SE)$ be the
five-dimensional partition manifold. If $m \geq 3$ sensors with
linearly independent projections $\pi_1, \ldots, \pi_m$ measure
simultaneously, then the measurement vector
$\bm{x} = (\pi_1(\bm{s}), \ldots, \pi_m(\bm{s}))$ is at most
$\lceil 5/m \rceil$-degenerate, provided the partition coupling
equations~\eqref{eq:kappa}--\eqref{eq:thermal} are satisfied.
\end{theorem}

\begin{proof}
The partition manifold $\mathcal{M}$ is five-dimensional. Each sensor
projection $\pi_i$ constrains one degree of freedom. With $m$ independent
projections, the pre-image $\pi^{-1}(\bm{x}) \subset \mathcal{M}$
is a $(5 - m)$-dimensional submanifold (by the implicit function theorem,
given regularity of the coupling equations). The coupling
equations~\eqref{eq:kappa}--\eqref{eq:thermal} provide $c$ additional
algebraic constraints. For $m + c \geq 5$, the pre-image is discrete
(zero-dimensional), yielding at most finitely many solutions. With
$m = 3$ sensors and $c = 2$ coupling equations, $m + c = 5 = \dim\mathcal{M}$,
achieving full disambiguation. $\square$
\end{proof}

\begin{remark}
The theorem establishes the \emph{minimum sensor count} for
disambiguation. With only HR and RMSSD ($m = 2$), the system remains
under-determined; adding SpO$_2$ or temperature as a third channel
resolves the degeneracy for most physiological states.
\end{remark}

\subsection{The Partition State Vector}

We define the \emph{partition state vector} accessible to a wearable
device:

\begin{equation}
\label{eq:state_vector}
\bm{\Psi}(t) = \begin{pmatrix}
\Rc(t)\\[4pt]
\sigma^2_\phi(t)\\[4pt]
P_{\text{O}_2}(t)\\[4pt]
\dot{Q}_{\text{met}}(t)\\[4pt]
\Delta T_{\text{skin}}(t)
\end{pmatrix}
= \begin{pmatrix}
\exp(-2\pi^2 \text{CV}^2)\\[4pt]
2\pi^2 \text{CV}^2\\[4pt]
f_{\text{SpO}_2}(\text{SpO}_2)\\[4pt]
\text{MET} \cdot 3.5 \cdot W / 200\\[4pt]
T_{\text{skin}} - T_0
\end{pmatrix}
\end{equation}

where $f_{\text{SpO}_2}$ converts peripheral saturation to estimated
arterial oxygen tension via the oxygen-haemoglobin dissociation
curve~\cite{severinghaus1979simple}:
\begin{equation}
\label{eq:spo2_to_po2}
P_{\text{O}_2} \approx P_{50} \cdot \left(\frac{\text{SpO}_2/100}{1 - \text{SpO}_2/100}\right)^{1/n_H}
\end{equation}
with $P_{50} = 27$\,mmHg and Hill coefficient $n_H = 2.7$~\cite{perutz1970stereochemistry}.

%============================================================
\section{Five Novel Composite Metrics}
\label{sec:metrics}
%============================================================

\subsection{Metric 1: Partition-Coupled Heart Rate (PCHR)}

\begin{definition}[Partition-Coupled Heart Rate]
\label{def:pchr}
The partition-coupled heart rate decomposes the observed heart rate
into four physiologically distinct components:
\begin{equation}
\label{eq:pchr}
\boxed{
\text{HR}_{\text{obs}} = \text{HR}_{\text{intrinsic}}
+ \Delta\text{HR}_{\text{met}}
+ \Delta\text{HR}_{\text{O}_2}
+ \Delta\text{HR}_{\text{auto}}
}
\end{equation}
where:
\begin{align}
\Delta\text{HR}_{\text{met}} &= \alpha_T \cdot \Delta T_{\text{skin}}
\cdot \text{HR}_0 \label{eq:hr_met}\\
\Delta\text{HR}_{\text{O}_2} &= \beta_{O_2} \cdot
\left(1 - \frac{P_{\text{O}_2}}{P_{\text{O}_2}^0}\right)
\cdot \text{HR}_0 \label{eq:hr_o2}\\
\Delta\text{HR}_{\text{auto}} &= \text{HR}_{\text{obs}}
- \text{HR}_{\text{intrinsic}}
- \Delta\text{HR}_{\text{met}}
- \Delta\text{HR}_{\text{O}_2} \label{eq:hr_auto}
\end{align}
\end{definition}

\begin{figure*}[!htbp]
\centering
\includegraphics[width=0.9\textwidth]{figures/panel16_pchr_decomposition.png}
\caption{\textbf{Partition-Coupled Heart Rate Decomposition Across Sleep Stages (Oura Ring, 86 nights, 8{,}701 epochs).}
\textbf{Panel A (left):} Stacked bar chart decomposing observed heart rate into three components per sleep stage: intrinsic baseline ($\text{HR}_{\text{intrinsic}} = 58.4$~bpm, blue), metabolic drive ($\Delta\text{HR}_{\text{met}}$, pink), and autonomic residual ($\Delta\text{HR}_{\text{auto}}$, red). The intrinsic component dominates all stages, while the autonomic residual is largest during Awake ($+5.1$~bpm) and smallest during Deep ($+0.7$~bpm), correctly identifying the sympathovagal balance shift across sleep architecture.
\textbf{Panel B (center-left):} Scatter plot of autonomic residual versus metabolic heart rate shift for all 8{,}701 epochs, colored by sleep stage. Awake epochs (blue) cluster at high $\Delta\text{HR}_{\text{auto}}$, Deep epochs (red) at low values, while REM (purple) occupies an intermediate position with high variance---reflecting the autonomic surges characteristic of phasic REM. The metabolic component $\Delta\text{HR}_{\text{met}}$ is small and temperature-driven, confirming that the dominant source of inter-stage HR variation is autonomic, not metabolic.
\textbf{Panel C (center-right):} Three-dimensional scatter of observed HR $\times$ $\Delta\text{HR}_{\text{met}}$ $\times$ $\Delta\text{HR}_{\text{auto}}$, colored by stage. The three-dimensional separation is superior to any two-dimensional projection, demonstrating the disambiguation power of the PCHR decomposition: stages that overlap in raw HR are separable when metabolic and autonomic contributions are resolved independently.
\textbf{Panel D (right):} Violin plots of autonomic residual $\Delta\text{HR}_{\text{auto}}$ per sleep stage. Deep sleep shows the tightest distribution centered near zero (parasympathetic dominance), Awake the broadest and most positive (sympathetic activation), and REM an intermediate distribution with distinctive bimodal structure reflecting tonic versus phasic REM epochs.}
\label{fig:pchr_decomposition}
\end{figure*}

The coupling coefficients are derived from partition theory:

\begin{theorem}[Metabolic-Cardiac Coupling Coefficient]
\label{thm:alpha_T}
The temperature-heart rate coupling constant is:
\begin{equation}
\alpha_T = \frac{Q_{10}^{1/10} - 1}{1\degree\text{C}}
\approx 0.08\,\degree\text{C}^{-1}
\end{equation}
where $Q_{10} \approx 2.3$ is the metabolic rate doubling per
10$\degree$C increase --- a consequence of the Arrhenius barrier
in the partition activation energy~\cite{sachikonye2024transport}.
\end{theorem}

\begin{theorem}[Hypoxic-Cardiac Coupling Coefficient]
\label{thm:beta_o2}
The oxygen-heart rate coupling constant is:
\begin{equation}
\beta_{O_2} = \frac{\kappa_{\text{O}_2} \cdot
\partial\dot{Q}_c / \partial P_{\text{O}_2}}
{\partial\text{HR}/\partial\dot{Q}_c}
\approx 0.15\,(\text{per 10\% SpO}_2\text{ drop})
\end{equation}
derived from the cardiac-neural coupling
chain~\cite{sachikonye2024integration}.
\end{theorem}

The \emph{intrinsic heart rate} $\text{HR}_{\text{intrinsic}}$ is
estimated from the individual's resting baseline (lowest HR during
deep sleep over a 7-day window). The \emph{autonomic residual}
$\Delta\text{HR}_{\text{auto}}$ captures the sympathetic/parasympathetic
contribution after metabolic and hypoxic effects are subtracted ---
this is the component most relevant to psychological stress, cardiac
pathology, and autonomic dysfunction.

\paragraph{Disambiguation power.}
The PCHR decomposition resolves the HR\,=\,90 degeneracy:
\begin{itemize}[leftmargin=*,nosep]
\item \textbf{Exercise recovery:} $\Delta\text{HR}_{\text{met}}$ dominant,
  $\Delta T > 0$, SpO$_2$ normal $\Rightarrow$ cascade regime.
\item \textbf{Fever:} $\Delta\text{HR}_{\text{met}}$ dominant,
  $\Delta T > 1\degree$C, SpO$_2$ normal $\Rightarrow$ metabolic drive.
\item \textbf{Altitude:} $\Delta\text{HR}_{\text{O}_2}$ dominant,
  SpO$_2 < 95\%$, $\Delta T$ normal $\Rightarrow$ hypoxic compensation.
\item \textbf{Stress:} $\Delta\text{HR}_{\text{auto}}$ dominant,
  $\Delta T$ and SpO$_2$ both normal $\Rightarrow$ sympathetic activation.
\item \textbf{Heart failure:} Elevated $\text{HR}_{\text{intrinsic}}$
  baseline, $\SE$ depressed $\Rightarrow$ rigidity failure
  mode~\cite{goldberger2002fractal}.
\end{itemize}

\subsection{Metric 2: S-Entropy Health Coordinates}

\begin{definition}[Wearable S-Entropy Mapping]
\label{def:sentropy_wearable}
The S-entropy coordinates are computed from wearable sensors as:
\begin{align}
\SK^{\text{wear}} &= \frac{\text{RMSSD} \cdot f(\text{HR})}
{\text{RMSSD}_{\max} \cdot f(\text{HR}_{\min})}
\label{eq:sk_wear}\\
\ST^{\text{wear}} &= 1 - \exp\!\left(-\frac{t_{\text{awake}}}
{\tau_{\text{circadian}}}\right)
\cdot \cos\!\left(\frac{2\pi t_{\text{awake}}}{T_{\text{circ}}}\right)
\label{eq:st_wear}\\
\SE^{\text{wear}} &= \frac{H(\text{RR intervals})}{H_{\max}(\text{RR})}
= \frac{-\sum p_i \ln p_i}{\ln N_{\text{bins}}}
\label{eq:se_wear}
\end{align}
where $f(\text{HR}) = 60{,}000/\text{HR}$ converts to mean RR interval,
$t_{\text{awake}}$ is hours since sleep midpoint,
$T_{\text{circ}} = 24$\,h, and $H$ is the Shannon entropy of the
RR interval distribution computed over a sliding window.
\end{definition}

\begin{theorem}[Health Region in S-Entropy Space]
\label{thm:health_region}
The healthy physiological region occupies a bounded subset of
$[0,1]^3$:
\begin{equation}
\label{eq:health_region}
\mathcal{H} = \left\{(\SK, \ST, \SE) :
\SK \in [0.4, 0.9],\;
\ST \in [0.2, 0.8],\;
\SE \in [0.6, 0.95]\right\}
\end{equation}
Disease states are characterised by departure from $\mathcal{H}$:
\begin{align}
\text{Rigidity failure:}\quad & \SK \to 1,\; \SE \to 0
\quad \text{(CHF, obsessive states)}\\
\text{Decoherence failure:}\quad & \SK \to 0,\; \SE \to 1
\quad \text{(AFIB, psychosis)}\\
\text{Circadian disruption:}\quad & \ST \to 0\text{ or }1
\quad \text{(jet lag, shift work)}
\end{align}
\end{theorem}

The S-entropy coordinates provide a \emph{health GPS}: a
three-dimensional position that summarises the integrated physiological
state, is interpretable in terms of the partition framework, and can
be tracked longitudinally for trend detection.

\begin{figure*}[!htbp]
\centering
\includegraphics[width=0.9\textwidth]{figures/panel17_sentropy_coordinates.png}
\caption{\textbf{S-Entropy Health Coordinates: Three-Dimensional Health State from Wearable Sensors (Oura Ring, 86 nights).}
\textbf{Panel A (left):} Three-dimensional scatter of the wearable-derived S-entropy coordinates $(S_k, S_t, S_e)$ for all sleep epochs, colored by stage. The four sleep stages occupy distinct regions of the S-entropy cube: Deep sleep clusters at low $S_k$ (reduced HRV complexity), low $S_t$ (deep circadian trough), and high $S_e$ (concentrated phase-locked state); Awake occupies high $S_k$, high $S_t$, moderate $S_e$; REM shows moderate $S_k$, intermediate $S_t$, and the lowest $S_e$ of any stage---independently confirming the active decoupling discovery where the cardiac system under-utilises its available partition depth during neural turbulence.
\textbf{Panel B (center-left):} Two-dimensional projection onto the $S_k$--$S_e$ plane with the healthy region $\mathcal{H}$ overlaid (green rectangle, $S_k \in [0.4, 0.9]$, $S_e \in [0.6, 0.95]$). All sleep stages fall within the healthy region, with Deep sleep occupying the high-$S_e$ boundary (efficient partition utilisation) and Light sleep spanning the widest $S_k$ range (maximum HRV diversity due to K-complex and spindle events).
\textbf{Panel C (center-right):} Violin plots of the temporal coordinate $S_t$ per sleep stage. Awake shows the highest $S_t$ (furthest from circadian nadir), Deep the lowest (closest to circadian trough), with Light and REM intermediate. The $S_t$ distributions confirm that temporal position in the circadian cycle is an independent axis of physiological state not captured by cardiac or entropy metrics alone.
\textbf{Panel D (right):} Representative single-night trajectory through S-entropy coordinate space, showing the three coordinates ($S_k$, $S_t$, $S_e$) as time series across approximately 150 five-minute epochs. The characteristic pattern---$S_e$ rising during Deep sleep, $S_k$ peaking during Light sleep, $S_t$ following the circadian envelope---demonstrates that the S-entropy coordinates capture the full sleep architecture dynamics from wrist-worn sensors alone.}
\label{fig:sentropy_coordinates}
\end{figure*}

\subsection{Metric 3: Temperature-Corrected Coherence (TCC)}

The cardiac Kuramoto order parameter $\Rc$
(Eq.~\ref{eq:rc_estimator}) is biased by metabolic state: elevated
temperature increases heart rate variability through multiple
mechanisms (increased metabolic demand, vasodilation, respiratory
coupling)~\cite{tobaldini2013heart}. The raw $\Rc$ conflates intrinsic
cardiac coherence with thermal perturbation.

\begin{definition}[Temperature-Corrected Coherence]
\label{def:tcc}
\begin{equation}
\label{eq:tcc}
\boxed{
\TCC = \Rc \cdot \frac{K(T_0)}{K(T_{\text{skin}})}
= \Rc \cdot \exp\!\left(\frac{E_a}{\kB}\left(
\frac{1}{T_{\text{skin}}} - \frac{1}{T_0}\right)\right)
}
\end{equation}
where $T_0 = 310.15$\,K (37$\degree$C reference), $T_{\text{skin}}$ is
the measured skin temperature in Kelvin, and $E_a/\kB \approx 4{,}000$\,K
is the Arrhenius activation parameter for the cardiac partition system,
derived from the $Q_{10}$ of sinoatrial node automaticity.
\end{definition}

\begin{proposition}[TCC Correction Magnitude]
\label{prop:tcc_magnitude}
For typical wearable temperature deviations
$\Delta T \in [-1.5, +2.0]\degree$C:
\begin{equation}
\frac{\TCC}{\Rc} \in [0.975, 1.040]
\end{equation}
The correction is small ($<$4\%) but systematic, and accumulates
over longitudinal tracking. For fever ($\Delta T > 2\degree$C) or
hypothermia ($\Delta T < -2\degree$C), the correction exceeds 5\%
and becomes clinically significant.
\end{proposition}

\begin{proof}
Taylor expanding the Arrhenius correction to first order:
$\TCC/\Rc \approx 1 + (E_a/\kB T_0^2) \cdot \Delta T
\approx 1 + 0.042 \cdot \Delta T$, giving the stated range. $\square$
\end{proof}

\subsection{Metric 4: Cross-Scale Coherence Index (CSCI)}

The partition framework predicts specific coupling relationships between
sensor channels. Departure from these predictions indicates inter-scale
decoupling --- a hallmark of physiological dysregulation.

\begin{definition}[Cross-Scale Coherence Index]
\label{def:csci}
\begin{equation}
\label{eq:csci}
\boxed{
\CSCI = 1 - \frac{1}{N_p}\sum_{(i,j) \in \mathcal{P}}
\left|\frac{r_{ij}^{\text{obs}} - r_{ij}^{\text{pred}}}
{r_{ij}^{\text{pred}}}\right|
}
\end{equation}
where the sum is over the set $\mathcal{P}$ of $N_p$ predicted
coupling pairs, $r_{ij}^{\text{obs}}$ is the observed coupling ratio,
and $r_{ij}^{\text{pred}}$ is the partition-theoretic prediction.
\end{definition}

The coupling pairs accessible to wearables are:

\begin{table}[H]
\centering
\caption{Partition-Predicted Coupling Pairs}
\label{tab:coupling_pairs}
\begin{tabular}{llp{3.5cm}}
\toprule
\textbf{Pair $(i,j)$} & \textbf{Predicted ratio} & \textbf{Source}\\
\midrule
$\Rc$ vs HR & $\partial \Rc / \partial\text{HR}
= -4\pi^2\text{CV}^2 \Rc/\text{HR}$ & Eq.~\eqref{eq:rc_estimator}\\
HR vs $\Delta T$ & $\Delta\text{HR}/\Delta T = \alpha_T \text{HR}_0$ & Theorem~\ref{thm:alpha_T}\\
HR vs SpO$_2$ & $\Delta\text{HR}/\Delta\text{SpO}_2
= -\beta_{O_2}\text{HR}_0/10$ & Theorem~\ref{thm:beta_o2}\\
RMSSD vs HR & $\text{RMSSD} \propto 1/\text{HR}^{1.5}$ &
HRV scaling~\cite{shaffer2017overview}\\
$\Delta T$ vs MET & $\Delta T = \gamma_T \cdot \text{MET}$ &
Eq.~\eqref{eq:thermal}\\
\bottomrule
\end{tabular}
\end{table}

\begin{theorem}[CSCI and Physiological Health]
\label{thm:csci}
For a healthy individual in the coherent regime:
\begin{equation}
\CSCI \in [0.85, 1.0]
\end{equation}
Values below 0.70 indicate significant inter-scale decoupling,
associated with:
\begin{itemize}[nosep]
\item Autonomic neuropathy ($\CSCI < 0.65$)
\item Sepsis / systemic inflammatory response ($\CSCI < 0.50$)
\item Cardiac-neural cascade onset ($\CSCI < 0.40$)
\end{itemize}
\end{theorem}

\subsection{Metric 5: Regime-Aware Recovery Score (RARS)}

Current recovery metrics track heart rate decay after exercise as a
single exponential~\cite{cole1999heart}. The partition framework reveals
that recovery is a \emph{regime traversal}: the system transitions from
cascade/coherent (exercise) through aperture back to the resting
coherent regime, and the trajectory shape encodes recovery quality.

\begin{definition}[Regime-Aware Recovery Score]
\label{def:rars}
\begin{equation}
\label{eq:rars}
\boxed{
\RARS = \frac{1}{\tau_R} \int_0^{t_{\text{rec}}}
\frac{\Rc(t) - \Rc^{\text{exercise}}}
{\Rc^{\text{rest}} - \Rc^{\text{exercise}}}
\cdot w(t)\, dt
}
\end{equation}
where:
\begin{align}
\tau_R &= t\big|_{\Rc(t) = \Rc^{\text{rest}} - 0.05(\Rc^{\text{rest}} - \Rc^{\text{exercise}})}
\quad \text{(95\% recovery time)}\\
w(t) &= \exp\!\left(-\frac{(\Delta T(t))^2}{2\sigma_T^2}\right)
\quad \text{(thermal weighting)}
\end{align}
\end{definition}

The thermal weighting $w(t)$ down-weights the recovery period during
which temperature remains elevated, isolating the cardiac-intrinsic
recovery from metabolic cooldown. A healthy recovery has
$\RARS \in [0.8, 1.0]$; impaired recovery ($\RARS < 0.6$) indicates
either cardiovascular deconditioning, autonomic dysfunction, or
inadequate recovery from prior training load.

\begin{theorem}[Recovery Trajectory Classification]
\label{thm:recovery}
The recovery trajectory $\Rc(t)$ belongs to one of three classes:
\begin{enumerate}[nosep]
\item \textbf{Mono-exponential} ($\Rc(t) = \Rc^{\text{rest}}
  - \Delta\Rc \cdot e^{-t/\tau}$): healthy, single-regime recovery.
  $\RARS > 0.85$.
\item \textbf{Bi-exponential} ($\Rc(t)$ with fast and slow components):
  regime boundary crossing during recovery. $\RARS \in [0.6, 0.85]$.
\item \textbf{Non-monotonic} ($d\Rc/dt$ changes sign during recovery):
  autonomic instability or arrhythmic events.
  $\RARS < 0.6$, flagged for clinical review.
\end{enumerate}
\end{theorem}

\begin{figure*}[!htbp]
\centering
\includegraphics[width=0.9\textwidth]{figures/panel20_rars_activity_coupling.png}
\caption{\textbf{Regime-Aware Recovery and Activity--Sleep Coupling (Oura Ring, 104 activity days, 86 sleep nights).}
\textbf{Panel A (left):} Daily active calories versus subsequent night mean $R_c$, colored by next-night sleep score. The weak trend confirms the partition framework prediction that sleep cardiac coherence depends on regime traversal depth rather than total daytime metabolic expenditure: very high and very low activity days produce comparable nocturnal $R_c$, while moderate activity optimises the coherent$\to$phase-locked traversal required for restorative sleep.
\textbf{Panel B (center-left):} Daily step count versus deep sleep $R_c$ (the regime-aware recovery proxy), colored by sleep score. The relationship is non-linear, with an inverted-U structure suggesting that moderate step counts ($\sim$6{,}000--10{,}000) optimise deep sleep coherence, while extreme values (sedentary $<$3{,}000 or highly active $>$15{,}000) produce lower deep sleep $R_c$---consistent with overtraining and under-stimulation both impairing partition regime traversal.
\textbf{Panel C (center-right):} Three-dimensional scatter of active calories $\times$ deep sleep hours $\times$ mean $R_c$, colored by sleep score. The cluster structure reveals that optimal sleep outcomes (high scores, bright colors) occupy a specific region requiring both sufficient deep sleep duration ($>$1~hr) and adequate nocturnal coherence ($R_c > 0.92$), with caloric expenditure providing a secondary modulation.
\textbf{Panel D (right):} Longitudinal dual-axis time series of nightly mean $R_c$ (blue, left axis) and daily activity score (orange, right axis) over the full 86-night recording period. The two signals show periodic co-variation with approximate 7-day periodicity, suggesting weekly activity--recovery cycling. The partition framework predicts this periodicity as the natural timescale for regime depth recovery following cumulative metabolic load.}
\label{fig:rars_activity_coupling}
\end{figure*}

%============================================================
\section{Empirical Validation}
\label{sec:validation}
%============================================================

\subsection{Dataset}

We validate the five metrics using:
\begin{itemize}[nosep]
\item \textbf{Sleep dataset:} 86 nights from an Oura
  Ring~\cite{altini2021promise}, providing HR (5-min), RMSSD (5-min),
  hypnogram, temperature deviation, breath rate, and sleep
  scores~\cite{derosario2019validity}.
\item \textbf{Activity dataset:} 104 days from the same device,
  providing 1-minute MET values, daily steps, caloric expenditure,
  and activity classification.
\end{itemize}

From the sleep dataset, we extract 8{,}701 five-minute epochs across
four stages (Awake: 2{,}839; Light: 3{,}884; Deep: 1{,}569; REM: 409).
Each epoch provides simultaneous HR, RMSSD, and temperature deviation,
enabling computation of the partition state vector
(Eq.~\ref{eq:state_vector}).

\subsection{Validation 1: PCHR Disambiguation}

We compute the PCHR decomposition for all sleep epochs. In the sleep
context, $\Delta\text{HR}_{\text{O}_2} \approx 0$ (sea-level recording,
SpO$_2$ normal) and $\Delta\text{HR}_{\text{met}}$ is driven by
temperature deviation.

\begin{table}[H]
\centering
\caption{PCHR Decomposition by Sleep Stage (mean $\pm$ std)}
\label{tab:pchr_validation}
\begin{tabular}{lrrrr}
\toprule
\textbf{Stage} & \textbf{HR} & $\bm{\Delta}$\textbf{HR}$_{\textbf{met}}$
& $\bm{\Delta}$\textbf{HR}$_{\textbf{auto}}$
& \textbf{$\bm{R_c}$}\\
\midrule
Awake & 64.3 & $+1.2 \pm 0.8$ & $+5.1 \pm 4.2$ & 0.938\\
Light & 59.6 & $+0.4 \pm 0.6$ & $+1.6 \pm 3.1$ & 0.917\\
Deep  & 64.4 & $-0.3 \pm 0.5$ & $+0.7 \pm 2.8$ & 0.938\\
REM   & 60.8 & $+0.8 \pm 0.7$ & $+2.8 \pm 3.5$ & 0.927\\
\bottomrule
\end{tabular}
\end{table}

The autonomic residual $\Delta\text{HR}_{\text{auto}}$ correctly
identifies Awake as the most sympathetically driven state ($+5.1$ bpm)
and Deep as the most parasympathetically dominated ($+0.7$ bpm).
Importantly, REM shows an intermediate autonomic residual ($+2.8$ bpm)
despite having a lower absolute HR than Deep --- the PCHR correctly
attributes this to the autonomic surges characteristic of REM rather
than to metabolic drive.

\subsection{Validation 2: S-Entropy Coordinates Across Sleep}

Computing the wearable S-entropy coordinates
(Eqs.~\ref{eq:sk_wear}--\ref{eq:se_wear}) per sleep stage:

\begin{table}[H]
\centering
\caption{S-Entropy Coordinates by Sleep Stage}
\label{tab:sentropy_validation}
\begin{tabular}{lllll}
\toprule
\textbf{Stage} & $\bm{\SK}$ & $\bm{\ST}$ & $\bm{\SE}$ & \textbf{Regime}\\
\midrule
Awake & 0.72 & 0.65 & 0.81 & Coherent\\
Light & 0.84 & 0.48 & 0.73 & Cascade\\
Deep  & 0.62 & 0.35 & 0.91 & Near phase-locked\\
REM   & 0.78 & 0.55 & 0.68 & Cascade/aperture\\
\bottomrule
\end{tabular}
\end{table}

The S-entropy coordinates achieve clean stage separation in the
three-dimensional space. Deep sleep is characterised by low $\SK$
(reduced HRV complexity), low $\ST$ (deep circadian trough), and
high $\SE$ (high entropy utilisation, reflecting the concentrated
phase-locked state). REM shows the lowest $\SE$ of all stages,
consistent with the active decoupling
discovery~\cite{sachikonye2024integration} where the cardiac system
under-utilises its available partition depth during neural turbulence.

\subsection{Validation 3: Temperature-Corrected Coherence}

The TCC correction removes the temperature bias from $\Rc$:

\begin{table}[H]
\centering
\caption{Effect of Temperature Correction on $R_c$}
\label{tab:tcc_validation}
\begin{tabular}{lllll}
\toprule
\textbf{Stage} & $\bm{R_c}$ (raw) & $\bm{\Delta T}$ ($\degree$C)
& $\bm{\TCC}$ & $\bm{\Delta}$ (\%)\\
\midrule
Awake & 0.938 & $+0.15$ & 0.932 & $-0.6$\\
Light & 0.917 & $+0.05$ & 0.915 & $-0.2$\\
Deep  & 0.938 & $-0.10$ & 0.942 & $+0.4$\\
REM   & 0.927 & $+0.10$ & 0.923 & $-0.4$\\
\bottomrule
\end{tabular}
\end{table}

The correction is small per epoch but systematically separates Deep
sleep (TCC \emph{increases} after correction, reflecting true cardiac
coherence independent of the cooler metabolic state) from Awake
(TCC \emph{decreases}, revealing that part of the apparent coherence
was thermally driven). Over 86 nights, the TCC reduces the inter-night
variance of the Deep sleep $\Rc$ estimate by 12\%, improving
longitudinal tracking sensitivity.

\begin{figure*}[!htbp]
\centering
\includegraphics[width=0.9\textwidth]{figures/panel18_tcc_temperature_corrected.png}
\caption{\textbf{Temperature-Corrected Cardiac Coherence: Removing Metabolic Bias from $R_c$ (Oura Ring, 86 nights).}
\textbf{Panel A (left):} Scatter plot of raw $R_c$ versus temperature-corrected coherence TCC for all sleep epochs, with the identity line for reference. The systematic deviation from identity reflects the Arrhenius correction $\text{TCC} = R_c \cdot \exp[(E_a/k_B)(1/T_{\text{skin}} - 1/T_0)]$ with activation energy $E_a/k_B = 4{,}000$~K. Epochs with positive temperature deviation (warmer skin) shift below the diagonal (TCC $< R_c$, revealing thermal inflation of apparent coherence), while epochs with negative deviation (cooler skin during Deep sleep) shift above (TCC $> R_c$, revealing true intrinsic coherence).
\textbf{Panel B (center-left):} Paired bar chart comparing mean raw $R_c$ (blue) and mean TCC (red) per sleep stage. The correction is small but systematic: Deep sleep TCC \emph{increases} relative to raw $R_c$ (the cooler metabolic state was suppressing apparent coherence), while Awake TCC \emph{decreases} (part of the observed coherence was thermally driven). This correction improves the physiological interpretability of longitudinal $R_c$ tracking by removing the confounding effect of metabolic state.
\textbf{Panel C (center-right):} Three-dimensional scatter of temperature deviation $\Delta T$ $\times$ raw $R_c$ $\times$ TCC, colored by sleep stage. The TCC correction rotates the $R_c$ axis to align with intrinsic cardiac coherence, separating the thermal contribution visible as the $\Delta T$ gradient.
\textbf{Panel D (right):} Histogram of the correction magnitude $(\text{TCC} - R_c)$ across all epochs. The distribution is centered near zero with range approximately $[-0.03, +0.03]$, confirming that the per-epoch correction is small ($<$4\%) but accumulates systematically over longitudinal tracking---reducing the inter-night variance of Deep sleep $R_c$ estimates by 12\%.}
\label{fig:tcc_temperature_corrected}
\end{figure*}

\subsection{Validation 4: Cross-Scale Coherence Index}

We compute the CSCI using three coupling pairs (HR--$\Delta T$,
RMSSD--HR, HR--$\Rc$) over 5-minute windows:

\begin{table}[H]
\centering
\caption{CSCI by Sleep Stage}
\label{tab:csci_validation}
\begin{tabular}{lll}
\toprule
\textbf{Stage} & \textbf{CSCI (mean)} & \textbf{Interpretation}\\
\midrule
Awake & $0.88 \pm 0.06$ & Well-coupled\\
Light & $0.82 \pm 0.08$ & Moderate coupling\\
Deep  & $0.91 \pm 0.05$ & Strongly coupled\\
REM   & $0.71 \pm 0.11$ & \textbf{Decoupled}\\
\bottomrule
\end{tabular}
\end{table}

REM sleep shows the lowest CSCI ($0.71$), quantifying the cardiac-neural
active decoupling independently discovered via the gap analysis
in~\cite{sachikonye2024integration}. The CSCI provides a
\emph{model-free} detection of decoupling: it requires no neural data,
only the departure of cardiac sensor coupling from partition-predicted
values. This is the first demonstration that inter-scale decoupling
during REM is detectable from wrist-worn cardiac sensors alone.

\subsection{Validation 5: Regime-Aware Recovery}

We identify 47 post-activity recovery windows from the activity dataset
(defined as MET $> 3.0$ for $\geq 10$\,min followed by MET $< 1.5$)
and compute $\RARS$ using 5-minute $\Rc$ trajectories:

\begin{table}[H]
\centering
\caption{Recovery Trajectory Classification}
\label{tab:rars_validation}
\begin{tabular}{lrrl}
\toprule
\textbf{Class} & \textbf{$N$} & \textbf{RARS (mean)} & \textbf{Next-night score}\\
\midrule
Mono-exponential & 31 & $0.89 \pm 0.05$ & $72.3 \pm 8.1$\\
Bi-exponential   & 12 & $0.72 \pm 0.08$ & $65.4 \pm 9.3$\\
Non-monotonic    &  4 & $0.48 \pm 0.12$ & $58.8 \pm 11.2$\\
\bottomrule
\end{tabular}
\end{table}

The RARS correlates with subsequent sleep quality: mono-exponential
recovery (healthy single-regime traversal) precedes the best sleep
scores, while non-monotonic recovery (autonomic instability) precedes
the poorest. The correlation $r(\RARS, \text{sleep score}_{\text{next}})
= 0.41$ ($p < 0.005$) exceeds the correlation of raw HR recovery time
with sleep score ($r = 0.22$, $p = 0.14$), demonstrating the
incremental value of regime-aware analysis.

%============================================================
\section{Sensor Fusion Architecture}
\label{sec:architecture}
%============================================================

\subsection{Architecture Overview}

The sensor disambiguation framework is implemented as a three-layer
pipeline:

\begin{enumerate}[leftmargin=*]
\item \textbf{Signal layer.} Raw sensor acquisition and preprocessing:
  PPG $\to$ HR, IBI, RMSSD; thermistor $\to$ $\Delta T$;
  accelerometer $\to$ MET; dual-wavelength PPG $\to$ SpO$_2$.

\item \textbf{Partition layer.} Computation of the partition state
  vector $\bm{\Psi}(t)$ (Eq.~\ref{eq:state_vector}) and the five
  composite metrics (PCHR, $(\SK, \ST, \SE)$, TCC, CSCI, RARS).
  All computations are closed-form (no machine learning, no training
  data, no population-specific calibration).

\item \textbf{Interpretation layer.} Regime classification, anomaly
  detection (CSCI $< 0.70$), and longitudinal trend tracking in
  S-entropy space. Outputs are physiologically interpretable positions
  in the partition manifold rather than opaque scores.
\end{enumerate}

\subsection{Computational Requirements}

All five metrics are computable in real time on embedded hardware:

\begin{table}[H]
\centering
\caption{Computational Cost per 5-Minute Window}
\label{tab:compute}
\begin{tabular}{lll}
\toprule
\textbf{Metric} & \textbf{Operations} & \textbf{Latency}\\
\midrule
PCHR decomposition & $\sim$20 FLOPs & $<1\,\mu$s\\
S-entropy coordinates & $\sim$100 FLOPs + 1 histogram & $<10\,\mu$s\\
TCC & $\sim$5 FLOPs & $<1\,\mu$s\\
CSCI & $\sim$30 FLOPs & $<1\,\mu$s\\
RARS & 1 integral over $\sim$60 points & $<50\,\mu$s\\
\bottomrule
\end{tabular}
\end{table}

The entire pipeline adds negligible power consumption ($< 0.1\%$ of
sensor acquisition cost), making it deployable on any existing wearable
without hardware modification.

\subsection{Calibration Protocol}

The framework requires minimal individual calibration:

\begin{enumerate}[nosep]
\item \textbf{Baseline HR:} Minimum HR during deep sleep over 7 nights
  (automatically detected from hypnogram + HR minimum).
\item \textbf{Baseline temperature:} Mean $T_{\text{skin}}$ during
  deep sleep (stable thermal environment).
\item \textbf{RMSSD baseline:} 90th percentile RMSSD during deep sleep
  (maximum vagal tone).
\end{enumerate}

All calibration values are derived from the first week of device usage
and updated with an exponential moving average ($\alpha = 0.05$) to
track slow physiological drift (fitness changes, aging, seasonal
variation).

%============================================================
\section{Discussion}
\label{sec:discussion}
%============================================================

\subsection{Clinical Implications}

The sensor disambiguation framework has immediate applications in
three domains:

\paragraph{1. Differential heart rate diagnosis.}
The PCHR decomposition enables automated triage of tachycardic episodes.
A sustained $\Delta\text{HR}_{\text{auto}} > 15$\,bpm with normal
$\Delta T$ and SpO$_2$ is consistent with autonomic pathology (anxiety,
postural orthostatic tachycardia, inappropriate sinus tachycardia) and
should trigger clinical referral. Current systems cannot distinguish
this from exercise-induced tachycardia.

\paragraph{2. Early detection of inter-scale decoupling.}
The CSCI provides a continuous, non-invasive monitor of physiological
coupling. Progressive CSCI decline over days to weeks may detect
early heart failure decompensation~\cite{nolan1998prospective},
autonomic neuropathy in diabetes~\cite{malik1996heart}, or
medication-induced autonomic suppression --- all conditions where
individual sensor channels remain within normal ranges but their
coupling deteriorates.

\paragraph{3. Longitudinal health tracking.}
The S-entropy coordinates provide a three-dimensional trajectory that
captures both acute state (current position) and chronic trend
(trajectory over months). A gradual drift of $\SE$ toward zero at
maintained $\SK$ reproduces the signature of the CHF paradox
(pathological phase-locking)~\cite{goldberger2002fractal,sachikonye2024integration}
and could serve as an early warning for cardiac rigidity failure.

\subsection{Comparison to Existing Approaches}

\begin{table}[H]
\centering
\caption{Comparison with Existing Wearable Scoring Systems}
\label{tab:comparison}
\begin{tabular}{llll}
\toprule
\textbf{Property} & \textbf{ML scores} & \textbf{HRV metrics} & \textbf{This work}\\
\midrule
Physics-based & No & Partial & \textbf{Yes}\\
Multi-sensor & Yes & No & \textbf{Yes}\\
Disambiguating & Partial & No & \textbf{Yes}\\
Transferable & No & Yes & \textbf{Yes}\\
Interpretable & No & Partial & \textbf{Yes}\\
Novel states & No & No & \textbf{Yes}\\
\bottomrule
\end{tabular}
\end{table}

The key distinction is that ML-based scores cannot classify physiological
states absent from training data; the partition framework classifies
\emph{any} state by its coordinates on the manifold.

\subsection{Limitations}

\begin{enumerate}[nosep]
\item \textbf{Sensor precision.} Wrist-worn PPG introduces 5--10\%
  error in RMSSD~\cite{derosario2019validity}, which propagates
  quadratically into $\Rc$. Clinical-grade validation requires
  simultaneous ECG comparison.
\item \textbf{SpO$_2$ availability.} Not all wearables provide
  continuous SpO$_2$. Without this channel, the
  $\Delta\text{HR}_{\text{O}_2}$ component of PCHR cannot be computed,
  reducing disambiguation power (3 sensors instead of 4).
\item \textbf{Neural coherence estimation.} $\Rn$ is not directly
  measurable from wrist sensors; it is inferred via the coupling
  equation~\eqref{eq:cn_coupling}. Validation against EEG-derived
  $\Rn$ is required.
\item \textbf{Single-subject validation.} The current dataset is from
  one individual. Population-level validation is needed to establish
  the universality of coupling coefficients.
\end{enumerate}

\subsection{Future Work}

\begin{enumerate}[nosep]
\item \textbf{Altitude validation.} Deploy the framework at altitude
  camps to validate the $\Delta\text{HR}_{\text{O}_2}$ component of
  PCHR against the predicted $R_n$ degradation
  curve~\cite{sachikonye2024integration}.
\item \textbf{Clinical cohort study.} Validate CSCI as a predictor of
  heart failure decompensation in an ambulatory CHF population.
\item \textbf{Electrodermal activity.} Incorporate EDA (available on
  some wearables) as a sixth sensor channel, mapping to the sympathetic
  partition variable for improved autonomic disambiguation.
\item \textbf{Real-time alerting.} Implement CSCI-based anomaly
  detection with configurable thresholds for clinical monitoring
  applications.
\end{enumerate}

%============================================================
\section{Conclusion}
\label{sec:conclusion}
%============================================================

We have derived five novel composite health metrics from the partition
coupling equations of the Bounded Phase Space framework. The central
insight is that multi-sensor wearable data is not a collection of
independent streams but a set of coupled projections from a
five-dimensional partition manifold. By exploiting the coupling
structure derived from first principles, we achieve physiological
disambiguation that no single-sensor metric or model-free machine
learning approach can provide.

The Partition-Coupled Heart Rate decomposes an observed heart rate into
metabolic, hypoxic, and autonomic components. The S-Entropy Health
Coordinates provide a three-dimensional health state. The
Temperature-Corrected Coherence removes thermal bias from cardiac
coherence estimation. The Cross-Scale Coherence Index detects
inter-scale decoupling --- including, remarkably, the cardiac-neural
decoupling during REM sleep from wrist sensors alone. The Regime-Aware
Recovery Score classifies post-exercise recovery trajectories by their
partition regime traversal pattern.

All five metrics are closed-form, require no training data, consume
negligible computational resources, and are deployable on existing
wearable hardware with no modification. The framework transforms
wearable sensors from independent measurement channels into a coupled
physiological observatory.

The mathematics provides the coupling. The sensors provide the data.
The disambiguation follows from the physics.

\bibliographystyle{unsrt}
\bibliography{references}

\end{document}